% ATTEMPT_STATUS: unresolved
% TOP_PROBLEM: {"problemId": "problem.grothendiecks-generalized-hodge-conjecture", "canonicalTitle": "Grothendiecks Generalized Hodge Conjecture", "releaseRank": 24, "primaryDomain": "number_theory_arithmetic_geometry", "primaryDomainLabel": "Number theory & arithmetic geometry", "displayStatus": "open_with_solved_subcases", "releaseStatus": "open_with_solved_subcases"}
% !TeX program = lualatex
\documentclass[11pt]{article}
\usepackage[margin=25mm]{geometry}
\usepackage{fontspec}
\setmainfont{Latin Modern Roman}
\usepackage{amsmath,amssymb,mathtools}
\usepackage{unicode-math}
\setmathfont{Latin Modern Math}
\usepackage{xurl,hyperref}
\hypersetup{colorlinks=true,urlcolor=blue}
\setcounter{secnumdepth}{0}
\setlength{\parindent}{0pt}
\setlength{\parskip}{5pt}
\setlength{\emergencystretch}{3em}
\begin{document}
\section*{\texorpdfstring{Grothendiecks Generalized Hodge Conjecture}{Grothendiecks Generalized Hodge Conjecture}}\label{24}
\subsection{Definitions and mathematical statement}
A rational Hodge substructure $V\subseteq H^k(X,{\mathbb{Q}})$ has Hodge coniveau at least $c$ if $V^{a,b}=0$ whenever $a<c$ or $b<c$. Let $M^cH^k$ be the sum of all such substructures. Define geometric coniveau by
\[
 N^cH^k(X,{\mathbb{Q}})=\sum_{\substack{Z\subseteq X\text{ closed algebraic}\\ \operatorname{codim}Z\ge c}}
 \ker\bigl(H^k(X,{\mathbb{Q}})\to H^k(X\setminus Z,{\mathbb{Q}})\bigr).
\]
The conjecture is $M^cH^k(X,{\mathbb{Q}})=N^cH^k(X,{\mathbb{Q}})$ for every smooth projective complex $X$ and all relevant $k,c$. Finite sums can be supported on the union of their supporting algebraic sets.
\par\smallskip\textit{Formulation and concept references:} \href{https://jomprob.org/index.php/jomp/article/view/Vol-1Issue-1Paper-2}{[S1]}.

\subsection{Short English statement}
Does every part of cohomology whose Hodge types allow high-codimension support actually have such algebraic support?
\subsection{Research attempt}

\paragraph{Scope and checked context (27 September 2026).}
This is the rational generalized Hodge conjecture for smooth projective
 complex varieties. The input is an entire rational Hodge substructure,
 not a list of individually chosen complex Hodge components. Voisin's
 survey [S1, Corollary 4.5 and Section 4.2; author PDF in E1] supplies the
 established inclusion $N^cH^k\subseteq M^cH^k$ and separates geometric
 realization from algebraicity of correspondences. The ordinary Hodge
 conjecture by itself is not invoked as a theorem. The additional
 catalogue references are not treated as resolutions of the general target.

The approach is to construct support by projecting an algebraic
 correspondence from an auxiliary variety. This gives a precise dimension
 test. In the first odd degree above the cycle-class case, the auxiliary
 Hodge structure can be realized by a curve, giving a conditional reduction
 with no unknown realization step. The remaining algebraicity step is
 isolated explicitly. Product examples then test both odd-degree support
 and a tempting but false operation on supporting sets.

\paragraph{Elementary support bookkeeping.}
Write $d=\dim X$. If classes $\alpha_j$ vanish on $X\setminus Z_j$, their
 sum vanishes on $X\setminus\bigcup_j Z_j$, by restriction. A finite union
 of codimension-at-least-$c$ algebraic sets still has codimension at least
 $c$. Since $H^k(X,\mathbb Q)$ is finite dimensional, a basis of any
 subspace contained in $N^cH^k$ uses only finitely many supports. Thus that
 entire subspace has one common algebraic support. There is no countable
 union or limiting algebraic set in this deduction.

For a smooth inclusion $j:Y\hookrightarrow X$ of codimension $c$, the
 Gysin map has the form
\[
 j_*:H^{k-2c}(Y,\mathbb Q)(-c)\longrightarrow H^k(X,\mathbb Q).
\]
It preserves Hodge types after the indicated twist: a type $(a,b)$ on
 $Y$ maps to type $(a+c,b+c)$ on $X$. Its image is supported on $Y$.
 Singular supports require the mixed-Hodge-theoretic support theorem,
 not a fictitious smooth Thom isomorphism on a singular variety.
 The required theorem is stated in [E1, Corollary 4.5]: rational coniveau
 classes can be obtained from smooth projective varieties of dimension
 $d-c$ by proper pushforward. This proves the known inclusion $N^c\subseteq
 M^c$, and also gives $N^cH^k=0$ when $k<2c$.

At $c=0$, both sides are the full cohomology. At $k=2c$, Hodge coniveau
 at least $c$ permits only type $(c,c)$. The support description then
 uses $H^0$ of the supporting resolutions, whose generators push forward
 to fundamental classes of codimension-$c$ components. Hence
\begin{equation}\label{ghc:even}
 M^cH^{2c}=H^{2c}(X,\mathbb Q)\cap H^{c,c}(X),\qquad
 N^cH^{2c}=\operatorname{span}_{\mathbb Q}\{[Z]:\operatorname{codim}Z=c\}.
\end{equation}
Thus this boundary degree is exactly the ordinary rational Hodge
 conjecture. Passing to greater degrees introduces cohomology \emph{on}
 the supports, not just their fundamental classes.

\paragraph{Lemma 24.1: a correspondence that constructs support.}
Let $Y$ be smooth projective of dimension $m$, and let
 $\Gamma\in\operatorname{CH}^{m+c}(Y\times X)\otimes\mathbb Q$.
 Its action
\[
 \Gamma_*\alpha=(\operatorname{pr}_X)_*
       (\operatorname{pr}_Y^*\alpha\smile[\Gamma])
\]
 takes $H^{k-2c}(Y,\mathbb Q)$ into $N^cH^k(X,\mathbb Q)$.

\emph{Proof.}
Choose a finite cycle representative $\Gamma=\sum_j a_j\Gamma_j$.
 Every component has dimension
 $m+d-(m+c)=d-c$. Its projection to $X$ is closed, because the
 projection is proper, and has dimension at most $d-c$. Set
 $Z=\bigcup_j\operatorname{pr}_X(\Gamma_j)$. On $Y\times(X\setminus Z)$
 the restricted cycle is zero, hence its cohomology class is zero.
 Compatibility of proper pushforward with restriction to this open set
 shows that $(\Gamma_*\alpha)|_{X\setminus Z}=0$. Therefore the whole
 image is supported on $Z$, and $\operatorname{codim}Z\geq c$.
 The cohomological degree is
 $(k-2c)+2(m+c)-2m=k$, as required.\hfill$\square$

The dimension of the \emph{cycle} is the key. An arbitrary Hodge-theoretic
 map with the same degree does not come equipped with components whose
 projections one can take. Also, a formal codimension-$c$ correspondence
 on $Y\times X$, instead of codimension $m+c$, has the wrong degree and
 generally the wrong support dimension. Neither error can be repaired by
 renaming a cohomology class a cycle.

There is also an unconditional top-odd-degree case for every $X$. For
 $d\geq2$, take a smooth complete-intersection curve $j:C\hookrightarrow X$.
 Weak Lefschetz makes $j^*:H^1(X,\mathbb Q)\to H^1(C,\mathbb Q)$
 injective. Under Poincar\'e duality its transpose is
 $j_*:H^1(C,\mathbb Q)\to H^{2d-1}(X,\mathbb Q)$, which is consequently
 surjective. Indeed, the transpose identity is
 $\int_X j_*\alpha\smile\beta=\int_C\alpha\smile j^*\beta$,
 and both pairings are nondegenerate. Every top-odd-degree class is
 therefore supported on this curve, in codimension $d-1$. The only
 Hodge types in that degree are $(d,d-1)$ and $(d-1,d)$, so
 $M^{d-1}H^{2d-1}=N^{d-1}H^{2d-1}=H^{2d-1}$. For $d=1$ the same
 conclusion is the already checked case $c=0$. This uses actual
 geometry of a hyperplane-section curve, rather than a hypothetical
 algebraic representative of a Hodge tensor.

There is a related projector trap. Even an algebraic correspondence
 $\Pi\in\operatorname{CH}^{d}(X\times X)\otimes\mathbb Q$ that acts as
 a projector onto a desired Hodge substructure has components of dimension
 $d$, and their projections can fill all of $X$. Its algebraicity alone
 gives support in codimension zero by the argument just proved. The
 diagonal is the basic example: it is an algebraic identity correspondence
 whose projection is all of $X$. To obtain codimension $c>0$, one needs
 an additional support theorem for the projector, or a factorization
 through cohomology of lower degree with a correspondence of the dimension
 in Lemma 24.1. A linear-algebraic splitting of cohomology therefore does
 not itself create the requested smaller support.

\paragraph{A conditional solution of the weight-one remainder.}
Fix $c\geq0$ with $2c+1\leq2d$, and let
 $V\subseteq H^{2c+1}(X,\mathbb Q)$ be a rational Hodge substructure
 of coniveau at least $c$. Then $W=V(c)$ has weight one and types
 $(1,0),(0,1)$ only. It is polarizable, since projective cohomology is
 polarizable and substructures inherit a polarization. Riemann's theorem,
 in the rational form recorded by Milne [E2, Theorem 6.8 and Corollary 6.9],
 realizes $W$ as $H^1(A,\mathbb Q)$ for a complex abelian variety $A$.
 The source states the covariant homology version; dualizing gives this
 cohomology version. The zero substructure needs no construction.

One can replace $A$ by a curve as a \emph{source}. If $\dim A=1$, take
 $C=A$. Otherwise choose a smooth complete-intersection ample curve
 $C\subset A$. Weak Lefschetz makes
 $H^1(A,\mathbb Q)\to H^1(C,\mathbb Q)$ injective. Its image is a
 polarizable Hodge substructure; the polarization-orthogonal complement
 gives a rational Hodge retraction. Consequently there is a surjective
 Hodge map $H^1(C,\mathbb Q)\to W$, and hence a map
\begin{equation}\label{ghc:phi}
 \phi:H^1(C,\mathbb Q)(-c)\longrightarrow H^{2c+1}(X,\mathbb Q)
 \quad\text{with image }V.
\end{equation}
This step uses no algebraicity assertion about that retraction.

Poincar\'e duality on $C$ identifies
 $H^1(C)^*\cong H^1(C)(1)$. Thus a rational Hodge morphism
 \eqref{ghc:phi} corresponds to a rational class
\[
 \gamma_\phi\in H^1(C,\mathbb Q)\otimes H^{2c+1}(X,\mathbb Q)
       \subset H^{2c+2}(C\times X,\mathbb Q)
\]
 of type $(c+1,c+1)$. Indeed the Hom tensor has its twist by $c+1$
 and is of type $(0,0)$; removing that twist gives the asserted type.
 The cup-product duality fixes the signs so that its correspondence action
 on $H^1(C)$ is exactly $\phi$. No Hodge projector is assumed algebraic
 to form this cohomological class.

\textbf{Proposition 24.2.}
If the ordinary rational Hodge conjecture in codimension $c+1$ holds
 for $C\times X$ for every smooth projective curve $C$, then
 $M^cH^{2c+1}(X,\mathbb Q)=N^cH^{2c+1}(X,\mathbb Q)$.

\emph{Proof.}
For the curve just constructed, the stated hypothesis represents
 $\gamma_\phi$ by a rational algebraic cycle
 $\Gamma\in\operatorname{CH}^{c+1}(C\times X)\otimes\mathbb Q$.
 Lemma 24.1 with $m=1$ puts its image $V$ inside $N^c$.
 This applies to every permitted $V$, hence to their sum $M^c$.
 The known opposite inclusion gives equality.\hfill$\square$

The unproved input is visibly an ordinary Hodge conjecture on a
 \emph{different variety}, in a specific codimension. Merely assuming
 that Hodge classes on $X$ itself are algebraic does not provide this
 cycle on $C\times X$. The weight-one realization removes one obstruction
 in this degree; it does not prove the ordinary Hodge conjecture needed
 for the correspondence. For larger $k-2c$, the realization by lower-degree
 geometric cohomology is itself an additional issue.

\paragraph{An unconditional odd-degree family with exact coniveau.}
Let $E$ be a complex elliptic curve and $X=E\times\mathbb P^r$.
 For $0\leq c\leq r$, K\"unneth gives
\begin{equation}\label{ghc:product}
 H^{2c+1}(X,\mathbb Q)=H^1(E,\mathbb Q)\otimes\mathbb Q h^c.
\end{equation}
The only types are $(c+1,c)$ and $(c,c+1)$, so the entire space has
 Hodge coniveau $c$. Include
 $j:E\times\mathbb P^{r-c}\hookrightarrow E\times\mathbb P^r$
 using a linear subspace. The Gysin map sends
 $\alpha\otimes1$ to $\alpha\otimes h^c$. It is surjective onto
 \eqref{ghc:product}, and explicitly supports every class in codimension
 $c$. Coniveau $c+1$ is impossible because $2c+1<2(c+1)$.
 Thus both filtrations have precisely the claimed step in this example.

No individual algebraic cycle has a fundamental class of odd degree,
 so a search for a codimension-$(c+1/2)$ cycle would miss this entire
 nonzero cohomology group. Its classes are supported on an algebraic
 set through that set's own $H^1$. Likewise, a nonzero rational line in
 $H^1(E)$ is not a Hodge substructure: complex conjugation forces equal
 dimensions of types $(1,0)$ and $(0,1)$, so a nonzero weight-one Hodge
 structure cannot have rational dimension one. The whole two-dimensional
 space is the correct substructure. This prevents an unjustified
 replacement of the rational Hodge input by a chosen rational basis line.

\paragraph{Stress test: intersections of supports do not support sums.}
Take $X=E_1\times E_2$ and divisors
 $D_1=E_1\times\{p_2\}$ and $D_2=\{p_1\}\times E_2$.
 Their classes vanish respectively away from $D_1,D_2$, so
 $[D_1]+[D_2]$ is supported on their union. It is nonzero since
\[
 D_1^2=D_2^2=0,\qquad D_1D_2=1,
 \qquad\int_X([D_1]+[D_2])^2=2.
\]
But $D_1\cap D_2$ is a point. By $N^2H^2=0$, no nonzero degree-two
 class is supported there. Replacing the union by an intersection would
 manufacture a false gain of one in coniveau. The finite-union argument
 at the beginning, not an intersection or a limit of shrinking sets,
 is the correct way to produce common support for a subspace.

\paragraph{Verification and the first remaining gap.}
The correspondence dimension, cohomological degree, and Tate-twist signs
 were checked separately. The curve reduction uses a surjection from
 $H^1(C)$; it does not claim that every abelian variety is a Jacobian.
 The algebraic cycle obtained conditionally has rational coefficients and
 may have several components with either sign. Effectivity and integral
 support classes are not required. In the product calculation $c=0$ gives
 support on all of $X$, while $c=r$ gives support on $E\times\{\mathrm{pt}\}$;
 both endpoints match the same Gysin formula.

The source theorem for arbitrary singular supports is explicitly retained
 as an external input. There is no deduction of it from a smooth normal
 bundle on a singular set. Likewise, no numerical period calculation is
 claimed to certify a rational Hodge substructure or an algebraic cycle.
 The proofs and countertests are symbolic and finite-dimensional.

Three focused continuations are to construct the class
 $\gamma_\phi$ algebraically for a specified family of $X$; to exploit
 geometric correspondences from curves when they are already available;
 and, in higher remaining weight, to establish geometric realization
 before demanding algebraicity of its map. In the conditional weight-one
 argument the first unproved statement is the algebraicity of
 $\gamma_\phi$ on $C\times X$. A formal Hodge morphism or a matrix of its
 periods does not supply that missing cycle.

\textbf{UNRESOLVED.} This attempt proves a support-producing correspondence
 lemma, the precise conditional odd-degree reduction, an explicit family
 with equality of both coniveau filtrations, and a counterexample to an
 invalid support operation. It does not establish the generalized Hodge
 conjecture for arbitrary varieties. Written by GPT-6 Astra Ultra using
 direct derivations and primary-source checks; no new general theorem or
 formal verification is claimed.

\subsection{Sources}
\begin{itemize}
\item[S1]Claire Voisin, Hodge and generalized Hodge conjectures, coniveau and algebraic cycles, Journal of Open Mathematical Problems 1(1) (2025), 16-51.
\url{https://jomprob.org/index.php/jomp/article/view/Vol-1Issue-1Paper-2}.
\textit{Formulation source.}
\item[S2]Incidence equivalence, a survey.
\url{https://arxiv.org/abs/2607.22233}.
\textit{Further reference; not a proof endorsement.}
\item[S3]On the Hodge and Tate conjectures for moduli spaces of curves.
\url{https://arxiv.org/html/2605.20453v1}.
\textit{Further reference; not a proof endorsement.}
\item[E1]Claire Voisin, \emph{Hodge and generalized Hodge conjectures, coniveau and algebraic cycles}, author-hosted full text of [S1], especially Proposition 2.11, Corollary 4.5, and Section 4.2.
\url{https://webusers.imj-prg.fr/~claire.voisin/Articlesweb/JOMP_V01N01P02Paper.pdf}.
\textit{Primary source for support and Hodge-substructure formalism; consulted 27 September 2026.}
\item[E2]J. S. Milne, \emph{Introduction to Shimura Varieties}, revised 16 September 2017, Theorem 6.8 and Corollary 6.9.
\url{https://jmilne.org/math/xnotes/svi.pdf}.
\textit{Primary author account of the realization of polarizable weight-one Hodge structures by abelian varieties, stated in the dual homology convention; consulted 27 September 2026.}
\end{itemize}
\end{document}
